\section{Limitation}

% Limited to one task
%Our findings are inherently limited by the settings of our empirical evaluation. 
The study focuses on seven LLMs and a single dataset, \textsc{NormAd-eti}, which, while comprehensive, does not represent all of the world's diverse cultural contexts and culturally-relevant prediction tasks. We selected \textsc{NormAd-eti} as our testbed for several reasons: \textbf{1)} \textsc{NormAd-eti} provides ground truth labels, enabling consistent and efficient evaluation. In contrast, other cultural benchmarks rely on aligning LLM outputs to responses from sociological surveys \cite{arora-etal-2023-probing, masoud2024culturalalignmentlargelanguage, kharchenko2024llmsrepresentvaluescultures, selfalignment, investigating} or probing-based methods \cite{culturegen, choenni-etal-2024-echoes, culturalbench}. \textbf{2)} \textsc{NormAd-eti} presents scenarios as stories, which align more closely with real-life contexts compared to cloze-style benchmarks \cite{culturegen}. \textbf{3)} \textsc{NormAd-eti} provides an extensive coverage, comprising of stories that reflect social and cultural norms from 75 countries. \textbf{4)} \textsc{NormAd-eti} is built on global community interviews with translators and rigorously validated by community experts, religious leaders, and academic researchers \cite{normad}. 
Despite these strengths, our findings should be interpreted with the limited scope of this task: \textsc{NormAd-eti} has 30-40 stories per country, which limits the generalization of the results, and is based on a ternary classification task  which does not account for decisions more nuanced than ``Yes'', ``No'', or ``Neither''.

% However, we will make the evaluation data, code, and framework publicly available so that new strategies and datasets can be seamlessly added.

% \mc{Need to give examples of these limitations, not only strengths since this is the limitations section :)}


% Variations of different debate formats
The scope of our multi-agent debate setup is as comprehensive as our computational budget allows, while we could not cover every possible variant of debate. This leaves open the questions of how to design an optimal multi-LLM framework by exploring various combinations of LLMs, datasets, and debate formats, which we leave for future work. Some interesting lines of future studies could be assigning specific roles to individual LLMs to represent particular countries or cultures \cite{culturepark}, or structuring debates where LLMs advocate for ``Yes'', ``No'', or ``Neither'' perspectives.

% \mc{Interesting, you could highlight some of these in conclusion as well}


% Computational inefficiency of multi-LLM collaboration
Our proposed multi-agent debate strategy involves prompting multiple LLM-based agents for feedback, which may introduce computational overhead, while the exact cost depends on the inference costs of debater LLM agents.\footnote{Detailed comparison of the computation and time efficiency of each method is in Appendix \ref{appendix:efficiency}.} We posit, however, that the value of multi-agent debate lies in leveraging diverse perspectives and reasoning paths across LLMs. Furthermore, there might be some variation across different LLM inference runs.

% Notably, our experiments demonstrate that two 7B LLMs can outperform \textsc{GPT-4} with significantly more parameters.